\section{Output-sensitive support scans inside the adaptive orbit kernel}
\label{sec:orbit-support}

The native cocycle stage makes remaining component-query costs visible.
A diagnostic profile of the 16-crossing stabilized circle finds 154,215
disjoint-support tests among 154,363 calls to the periodic merger. Such
tests cannot succeed. The previous adaptive scheduler avoids repeated
cubic scans by switching to a quadratic candidate queue, but still tests
every pair on a no-merger pass. The new scheduler adds an interval-specific
support scan while retaining that queue as its dense-case fallback.

\subsection{Pruning failures without changing the proof trace}

A periodic pairing has support $[a,d]$ and positive period $p=c-a$.
Both maintained merger rules require positive support overlap: the sharp
rule requires at least $p+q-\gcd(p,q)\ge1$ points, and the classical rule
requires at least $p+q$. Therefore deleting disjoint pairs from a candidate
scan cannot remove a successful merger under either rule.

As before, find the first two periodic rows lazily and try that pair before
building any index. Immediate first-pair merger families retain linear
eligibility work and construct no support index. After a failure, collect
all periodic rows and sort them by support start, with original row index
as a tie-breaker. Sweep from left to right, retaining earlier rows in an
active set and a heap ordered by support end. Remove a row only when its
end is strictly less than the next start: supports meeting at one integer
point must remain eligible.

Each active row overlaps the new row, and every overlapping pair appears
exactly once, when its later sweep endpoint is visited. Store its two
original indices in increasing order and sort the candidate list
lexicographically. Resume the merger scan in that order, skipping the
already failed first pair. Until a merger succeeds, row values and list
indices are fixed. Omitted pairs all fail, so the first success is exactly
the old greedy success. Its replacement, survivor position and trace event
are unchanged. Induction after each restart proves identical closure output
and successful events. The resulting complete AHT certificates and cycle
allowances are preserved; scheduling diagnostics and work-guard counts may
change. The independent replay checkers are unchanged.

For $k$ input rows, give the closure allowance $A=\binom{k}{2}$. Charge
each initial first-pair test and \emph{every enumerated overlap}, including
ones later discarded because an earlier candidate merges. If another charge
would exceed $A$, abandon that unfinished index and use the existing
generation-checked queue on the current rows. No merger from a partial index
has been applied. The fallback therefore starts from precisely the same
greedy state and cannot change a verdict or trace. Charging index construction,
rather than only attempted mergers, prevents repeated dense index builds
from escaping the allowance.

A no-merger support pass with $E$ overlaps costs
$O(k\log(k+1)+E\log(k+1))$ scalar operations and $O(k+E)$ storage.
For disjoint supports $E=0$, it makes one merger call and takes
$O(k\log(k+1))$ operations rather than testing $\binom{k}{2}$ pairs.
There are at most $k$ support builds before closure terminates, and their
total charged overlap count is at most $A$. With the fallback, the
conservative bound remains $O(k^2\log(k+1))$ scheduler operations,
$A+\max(k-1,0)^2$ merger calls and $O(k^2)$ memory. Binary endpoint
comparisons and the exact gcd merger have their ordinary bit costs. This
is a bound for one closure; it does not count all AHT cycles or prove a
general quasi-polynomial unknot recognizer.

There is also an explicit family on which this changes the cost of the
whole orbit computation. On $[0,2k)$ take the $k$ unit pairings
$[2i,2i]\leftrightarrow[2i+1,2i+1]$. Each cycle truncates the rightmost
pair and the next contraction emits its remaining singleton. The algorithm
has $k+1$ cycles, including the final empty-pairing contraction. The previous
no-merger scans make
\[
 \sum_{j=1}^k\binom{j}{2}=\binom{k+1}{3}
\]
merger calls; the support scheduler makes exactly $k-1$. Remaining coverage,
transmission and independent replay scans give a conservative
$O(k^2\log(k+1))$ scalar-operation bound for the new complete query and
replay, compared with the old cubic merger work. The proof trace is identical.
This is a restricted interval-system theorem, not a claim that arbitrary
normal-surface systems have disjoint supports.

\subsection{One static-gap scan and an identity contraction}

The original producer separately scanned domains and ranges to record
certificate gaps, then repeated the sort and coverage scan to perform the
contraction. It also rebuilt every immutable pairing even when there was
no static gap. The new routine obtains occupied blocks and their complement
in one sorted scan and returns the same gap list for the trace event.

If the complement is empty, the order-preserving contraction is the identity,
so it reuses the immutable pairing rows. Otherwise a prefix map on occupied
blocks gives every endpoint's new rank and reconstructs the pairings. All
gaps contain only singleton orbits; their total size is added to the orbit
count. This is the same mathematical reduction and the same certificate
event as before, with fewer sorts and allocations. Input rows are immutable
and the caller already owns its working list; no caller-supplied pairing
is mutated. Cooperative cancellation covers both the coverage scan and
the nontrivial endpoint map.

\subsection{Source preservation, bounded queries and independent checks}

Five new focused tests pass in 0.158 seconds and all 1,171 maintained tests
pass in 210.854 seconds. The new controls compare 400 random closures under
both merger rules with the literal restart scheduler, including exact trace
order and the per-closure candidate bound. They compare 250 contractions
with explicit finite order maps, retain one-point support intersections,
exercise the dense fallback and cancellation, and use $2^{20000}$-scaled
endpoints without expanding their represented points. Existing integration
tests retain exact weighted and normal-component certificates, both merger
rules, false-valued callbacks and all relevant cycle allowances.

The pinned audit performs 600 complete orbit comparisons on 300 additional
random systems under both rules. Counts, cycles and entire certificates
match the published baseline, and every result replays independently.
Small random systems do not exhibit the main support-sparsity effect:
their total merger calls only fall from 2,625 to 2,618. The maintained
finite-graph tests separately compare answers with literal connectivity.

The source audit reruns all 84 previous diagrams with the same shared
two-million-guard allowance. Seventeen positive results replay independently
and match the earlier source-certificate hashes. Sixty-three cases finish
the restricted candidate queries without a disc and four remain capped,
compared with sixty-two completed misses and five caps previously.
The newly finished case is \code{survivor-01}: it uses 1,812,276 guards and
finds no disc in either the 331-piece raw surface or the 257-piece optimized
surface. Its final source status remains inconclusive. No recognition
coverage increase is claimed.

An independently reproducible diagnostic wraps the merger to classify every
call on the raw 16-crossing stabilized circle. Disjoint calls fall from
154,215 to 467, and overlapping calls remain exactly 148. Both complete
source proofs have the same hash. The shared guarded work falls from
1,126,713 to 988,196, a much smaller ratio than merger-call counts because
the rest of discovery, transport and positive replay still runs. The
diagnostic and source audit overlap regression tests, so their elapsed
times are not used as performance comparisons.

\subsection{Serial complete-query measurements}

The final experiments run serially after tests, audits, diagnostics and the
runtime commit. Each has five shuffled paired rounds, old and new A/A
controls, and one excluded warmup per arm and input. All 200 measured calls
and forty warmups complete. Within each input, all old and new complete
certificate hashes agree across repetitions. The tables separate individual
arm medians from medians of paired ratios.

The first endpoint counts the complete interval system, constructs its
certificate, independently replays that certificate and computes its hash.
The disjoint cases use the family proved above. The five-cycle cases retain
the earlier merger stress family, once with small and once with 4096-bit
endpoints. No represented integer point is expanded.

\begin{center}
\small
\begin{tabular}{lrrrrr}
\toprule
Input & Old (ms) & New (ms) & Old/new & Old A/A & New A/A \\
\midrule
disjoint-16 & 2.003 & 1.368 & 1.491 & 1.005 & 1.001 \\
disjoint-64 & 58.093 & 14.533 & 4.128 & 0.991 & 1.026 \\
disjoint-256 & 3253.678 & 211.811 & 15.334 & 1.012 & 0.999 \\
Five-cycle, small & 3.028 & 2.604 & 1.147 & 1.006 & 1.017 \\
Five-cycle, 4096 bits & 10.133 & 7.785 & 1.296 & 0.999 & 1.004 \\
\bottomrule
\end{tabular}

\end{center}

All one hundred measured queries and twenty warmups finish; driver elapsed
time is 42.799 seconds. The disjoint cases improve by $1.491$, $4.128$ and
$15.334$. At $k=256$, old/new median times are 3.253678 and 0.211811 seconds,
and merger calls are exactly 2,796,160 and 255, as the formulas predict.
The five-cycle cases improve by $1.147$ and $1.296$. Old A/A ratios lie in
$0.991$--$1.012$ and new A/A in $0.999$--$1.026$. These complete-query
measurements include verifier cost; they are not isolated scheduler timings.

The second endpoint is the current full recognizer, with earlier shortcuts
disabled and the native stage enabled. Only the orbit producer is swapped
between releases. Both arms have no local native work cap, while retaining
the same twelve-second global deadline; every measured call completes below
that deadline. The three stabilized circles return raw normal-disc proofs,
the optimized-positive input needs its second candidate, and the trefoil
continues to the exact Khovanov fallback. Common result bookkeeping includes
fingerprinting the positive certificate.

\begin{center}
\small
\begin{tabular}{lrrrrr}
\toprule
Input & Old (ms) & New (ms) & Old/new & Old A/A & New A/A \\
\midrule
circle-8 & 358.809 & 252.720 & 1.430 & 0.995 & 1.000 \\
circle-16 & 1360.820 & 798.522 & 1.704 & 0.999 & 0.995 \\
circle-32 & 5902.774 & 2803.070 & 2.108 & 0.998 & 0.983 \\
Optimized positive & 285.147 & 215.253 & 1.334 & 1.004 & 1.003 \\
Trefoil & 84.994 & 73.013 & 1.169 & 0.996 & 1.001 \\
\bottomrule
\end{tabular}

\end{center}

All one hundred measured recognitions and twenty warmups finish; driver
elapsed time is 146.078 seconds. Paired gains range from $1.169$ to $2.108$,
outside old A/A variation $0.995$--$1.004$ and new A/A variation
$0.983$--$1.003$. The 32-crossing circle's median falls from 5.902774 to
2.803070 seconds, including native positive replay. The five-crossing
optimized success falls from 285.147 to 215.253 milliseconds. These are
whole-recognition gains on the explicitly exercised native-stage workload.
The 32-crossing source still exhausts the default two-million-guard allowance;
its unlimited-local-budget timing is not a new default-budget success.
No ordinary-portfolio gain or general recognition complexity bound is
inferred from this selected workload.


Reproduction uses \path{../fast/normal_orbit_research/supports.py}.

Its modes are \code{audit}, \code{orbits} and \code{recognize}. The call-category
diagnostic is \path{data/orbit_support_profile.py}. All 473 runtime, driver,
fixture and prior-audit pins are checked across the final experiments.
The baseline is commit \code{7b0746a27}, before either the support scheduler
or contraction change. The external \code{openai/math} head is unchanged
at \code{fd4aeeb2e}, and no new unknot
delivery arrived during this integration. The generic restart, hybrid and
queue schedulers remain available for comparisons; the specialized support
scheduler is the default inside the maintained orbit producer.

